

\newcommand{\diam}[1]{\textbf{Diam}(#1)}
\newcommand{\pred}[1]{\text{pre}_{D}(#1)}
\begin{definition}
  We define the diameter function $\textbf{Diam} \colon \mathbb{F}_{2}^{\mathcal{F}} \rightarrow \mathbb{R}$ to act on a finite chain $z \in \mathbb{F}_{2}^{\mathcal{F}}$ as:
  \begin{equation*}
    \begin{split}
      \diam{z } \colonequals \max_{v \in z}\max_{g \in S} d L_{g}(v) - \min_{v \in z}\sum_{g \in S} L_{g}(v) 
    \end{split}
  \end{equation*}
In addition, for a vertex $v$ in a space-time graph, we define $\pred{v}$ to be the vertex in $\pre{v}$ that achieves the maximal value of $\max_{g\in S}L_{g}$. If there are many, select one arbitrarily.
\end{definition}
The reduced form of an error can now be redefined to match the diameter function rather than the size. The following is an adaptation of \Cref{claim:reducederr} for the diameter function; the proof is identical to the original version.
\begin{claim}
  \label{claim:rereduced}
  Let $x \in \mathbb{F}_{2}^{\mathcal{F}_{d^{\prime}}}$ be a $d^{\prime}$-dimensional finite chain, consider the stabilizer  $z \in \mathcal{C}_{Z}^{\perp}$ which adding it to $x$ minimize the size , Namely: 
  \begin{equation*}
    \begin{split}
      z =  \arg \min_{z \in \mathcal{C}_{Z}^{\perp}} \diam{ x + z }  
    \end{split}
  \end{equation*}
  Then: 
  \begin{equation*}
    \begin{split}
      \diam{ x + z } = \diam{ \partial^{-}( x + z) } 
    \end{split}
  \end{equation*}
\end{claim}
The diameter function fits well within the definition of space-time graphs. The following is an adaptation of \Cref{claim:prevslices} to relate between the diameter of a slice and its predecessors.
\begin{claim} 
    \label{claim:prevslicesdiam}
    Consider a time graph $(G,T)$, Let $K$ be a slice in $G$, and let $K$ be a slice of $K$ such that $K \neq \History{K}$. Select vertices $x, y \in K$ that achieve the maximal values of $\max_{g}L_{g}$and $L_{d+1}$.  Then for some integer $k$ there exist slices $K_{0} , ..,K_{k}$ and forks $F_{1},..,F_{k}$ such that:
    \begin{enumerate}
      \item  The vertices $ \predd{x}, \pred{y} $ belong to $K_{0},..,K_{k}$.  
      \item Introduce an edge between $K_{i}$ and $K_{j}$ if and only if one of the forks $F_{1},..,F_{k}$ consists of a pole of $K_{i}$ and a pole of $K_{j}$. Then the graph formed from $K_{0},..,K_{k}$ is connected.  
      \item $\sum_{i=0}^{k}\diam{K_{i}} +  \sum^{k}_{i=1}\diam{F_{i}} \ge \max_{g\in S}L_{g}( \pred{y} ) +  L_{d+1}( \predd{x} )  $
    \end{enumerate}
  \end{claim}
  \begin{proof}
The proof of the first two conditions is identical to \Cref{claim:prevslices}, where $K_{1}, K_{2}, \ldots$ are chosen such that the connected graph is a spanning tree. For the third, first we use the fact that in the formed spanning trees there are (at most) two leaves and therefore, the degree is lower than two, and the graph is a simple path. Denote by $J_{0}, \ldots, J_{2k}$ the enumerated union of $K_{0}, \ldots, K_{k}, F_{1}, \ldots, F_{k}$. Now consider the $2k+1$ vertices $z_{0}^{1}, z_{0}^{2}, z_{1}^{1}, z_{1}^{2}, \ldots, z_{2k}^{1}, z_{2k}^{2}$ such that $z_{i}^{j} \in J_{i}$ and that $J_{i}$'s are ordered such that $J_{i} \cap J_{i+1} \neq \emptyset$. In addition, let us assume that $\predd{x} \in J_{0}, \pred{y} \in J_{2k}$. Observe that for any such collection of vertices we have that:
    \begin{equation*}
      \begin{split}
        \sum_{i=0}^{k}\diam{K_{i}} & +  \sum^{k}_{i=1}\diam{F_{i}}  \ge \sum_{i=0}^{2k} \max_{g\in S}d L_{g}( z_{i}^{1} ) +  L_{d+1}( z_{i}^{2} )  \\
      & \ge  \max_{g\in S}d L_{g}( z_{0}^{1} ) +  L_{d+1}( z_{0}^{2} )  +  \left( \sum_{i=1}^{2k-1} - L_{d+1}( z_{i}^{1} ) + L_{d+1}( z_{i}^{2} )   \right)  
      \end{split}
    \end{equation*}
    Now, since $ J_{i} \cap J_{i+1} \neq \emptyset $, we can choose $ z_{i}^{1} = z_{i-1}^{2} $ for any $ i \in [1,2k-1] $, so the above becomes a telescoping sum. Moreover, setting $ z_{0}^{1} = x $ and $ z_{2k}^{2} = y $ yields the desired inequality.
  \end{proof}




\begin{claim}
  \label{claim:diamexp}
  Consider a $d^{\prime}$-dimensional chain $z$ such both $z$ and $\phi(z)$ are finite, then it holds that: 
      \begin{equation*}
        \begin{split}
          \bigl| \diam{ z  } - \diam{ \phi(z)  } \bigr| < 2d + d^{\prime}
        \end{split}
      \end{equation*}
\end{claim}
\begin{proof}
  We start by showing that: 
\begin{equation*}
  \begin{split}
    \max_{v \in \phi(z)}L_{d+1}(v) = \max_{v \in z }L_{d+1}(v) - d^{\prime}
  \end{split}
\end{equation*}
Let us consider the action of $\phi$ on a rooted chain. Let $y$ be a rooted chain at vertex $v$, namely there are faces $A = \{ \rotfvs{} \}$, such that:
\begin{equation*}
  \begin{split}
    y = \sum_{ \rotfvs{} \in A}\rotfvs{}
  \end{split}
\end{equation*}
Thus, 
\begin{equation*}
  \begin{split}
    \phi(y) = \sum_{ \rotfvs{} \in A}\rotfv{S^{\prime}v}{S/S^{\prime}}
  \end{split}
\end{equation*}
So if $v$ achieves the maximal $L_{d+1}$-value in $z$, then the vertices that get the maximal value in $\phi(z)$ are of the form $S^{\prime}v$, for some $d^{\prime}$-size subset of generators $S^{\prime} \subset S$. Therefore, the following equality holds:
\begin{equation*}
  \begin{split}
    \max_{v \in \phi(z)}L_{d+1}(v) = \max_{v \in z }L_{d+1}(v) - d^{\prime}
  \end{split}
\end{equation*}
So it's left to show that for any generator $g \in S$ it holds that: 
\begin{equation*}
  \begin{split}
     \bigl| \max_{v \in \phi(z)}L_{g}(v) -  \max_{v \in z }L_{g}(v) \bigr| \le 1
  \end{split}
\end{equation*}
Let $v$ be a vertex that is a root of a face in $\phi(z)$. We argue that $v$ is also in $z$. Consider a face $\rotfvs{}$ in $\phi(z)$, then $\phi^{-1}(\rotfv{v}{S^{\prime}}) = \rotfv{(S/S^{\prime})^{-1} v}{S/S^{\prime}}$ is in $z$, and therefore $(S/S^{\prime})(S/S^{\prime})^{-1} v = v$ is in $z$. By the same arguments, we also have that if $v$ is a root of a face in $z$, then the vertex $S v$ is in $\phi(z)$.


Now, fix a generator $g \in S$ and let $u$ and $w$ be the vertices in $z$ and $\phi(z)$ that achieve the maximum for $L_{g}$. Denote by $u_{r}$ and $w_{r}$ the vertices which are the roots of the faces containing $u$ and $w$, and observe that $L_{g}(w) \le L_{g}(w_{r})+1$ for any $g \in S$. We have that:
    \begin{equation*}
      \begin{split}
        \max_{v \in z} L_{i}(v) &  \le   L_{g}(S u_{r}) \le  \max_{v \in \phi(z)} L_{g}(v)  \\
        \max_{v \in \phi(z)} L_{g}(v) & \le  L_{g}(w_{r})  + 1 \le  \max_{v \in z} L_{g}(v) + 1 
      \end{split}
    \end{equation*}
So, the inequality holds.
\end{proof}





  Let $p$ be a point in $A \cap B$. Let $X(t): [0,1] \rightarrow \mathbb{R}^{d}$ be the segment connecting $p$ to $\psi(p)$, namely:
\begin{equation*}
  \begin{split}
    X(t) \colonequals t p + (1-p) \psi (p)
  \end{split}
\end{equation*}
We argue that:
\begin{equation*}
  \begin{split}
    \Size{ A \cup \psi(B)} = \Size{ A \cup \psi(B) \cup X } 
  \end{split}
\end{equation*}
To see it, first consider the value of $L_{d+1}$ on the line $X$, for any $t \in [0,1]$, we have that: 
\begin{equation*}
  \begin{split}
  L_{d+1} (X(t)) = t L_{d+1}(p) + (1-p) L_{d+1}(\psi (p)) = L_{d+1}(p)
  \end{split}
\end{equation*}
So, the line $X$ does not introduce any new value of $L_{d+1}$. In addition, any other potential wall $L_{i}$ is a monotone convex function; therefore, it gets its maximal values on the boundary of $X$ (which is a compact, convex set), namely on either $p$ or $\psi(p)$. Hence, the maximal values of the potential wall of $A\cap \psi(B)$ and $A \cup \psi(B) \cup X$ are the same, as well as their sizes.

Suppose that there is a potential wall $L_{i}$ that achieve on
  Then:
  \begin{equation*}
    \begin{split}
      \Size{A\cup \psi(B)} &= \sum_{i \in [d+1]}\max_{v \in A\cup \psi(B)} L_{i}(v) \\
      & =  \sum_{ i \in [d+1]} \max \bigl\{ \max_{v \in A } L_{i}(v) , \max_{v \in \psi(B) } L_{i}(v) \bigr\}  \\
      & =  \sum_{ i \in [d+1]}  \max_{v \in A } L_{i}(v) + \max_{v \in \psi(B) } L_{i}(v) -  \min \{   \max_{v \in A } L_{i}(v), \max_{v \in \psi(B) } L_{i}(v) \} \\
      & \le  \sum_{ i \in [d+1]}  \max_{v \in A } L_{i}(v) + \max_{v \in \psi(B) } L_{i}(v) -   \overbrace{\min \{  L_{i}(p),  L_{i}(\psi(p)) \} }^{ \star  } 
        \end{split}
      \end{equation*}
Observe that the contribution of terms in $(\star)$ to the sum is negative. This is because:
\begin{equation*}
  \begin{split}
    \sum_{i \in [d+1]}\min \{  L_{i}(p),  L_{i}(\psi(p)) \} = \sum_{i \in [d]}\min \{  L_{i}(p),  L_{i}(\psi(p)) \} - \min \big\{ \sum_{i \in [d]} L_{i}(p),\sum_{i \in [d]} L_{i}(\psi(p))\big\} \le 0
  \end{split}
\end{equation*}
Thus,  
      \begin{equation*}
        \begin{split}
      \Size{A\cup \psi(B)} & \le  \sum_{ i \in [d+1]}  \max_{v \in A } L_{i}(v) + \max_{v \in \psi(B) } L_{i}(v)  \\
      & \le \Size{A} + \Size{\psi(B)}  =  \Size{A} + \Size{B} 
    \end{split}
  \end{equation*}

